%% Benne basculante à vérin — mêmes pièces, coloriées par classe d'équivalence.
\documentclass[tikz,border=4pt]{standalone}
\usepackage{liaisons}
%% Dessin d'ensemble simplifié de la benne basculante à vérin (coupe dans le plan (x, y)).
%% Inclus par mecanisme-benne-verin-dessin.tex et mecanisme-benne-verin-classes.tex, qui définissent
%% \piece{N}{chemin} (pièce coupée : hachurée ou colorée), \piecevue{N}{chemin} (pièce non coupée :
%% tige, axes, roues) et \bnlabels (les étiquettes propres à chaque vue).
%% Cotes en cm. Points : C = (5,6 ; 1,2) axe de basculement ; O = (4,2 ; 0,5) pied du vérin ;
%% P = (2,4 ; 1,5) attache de la tige sous la benne. Axe du vérin : u = (O -> P), 150,945°.
\newcommand\benne{%
  \begin{scope}[line width=0.6pt, line join=round]
    % --- sol
    \draw[line width=0.5pt, black!60] (0.05,0) -- (6.55,0);
    % --- 0 : châssis du camion (longeron coupé, cabine, chapes) ; camion à l'arrêt
    \piece{0}{(0.30,0.18) rectangle (6.25,0.50)}                              % longeron (coupé)
    \piece{0}{(0.30,0.50) -- (0.95,0.50) -- (0.95,2.15) -- (0.55,2.15) -- (0.30,1.45) -- cycle}  % cabine
    \piece{0}{(5.46,0.50) rectangle (5.74,1.20)}                              % chape de basculement (sur le châssis)
    % --- 1 : benne (caisse coupée) + chape d'articulation arrière
    \piece{1}{(1.0,2.60) -- (1.0,1.50) -- (5.6,1.50) -- (5.6,2.30) -- (5.42,2.30) -- (5.42,1.68) -- (1.18,1.68) -- (1.18,2.60) -- cycle}
    \piece{1}{(5.42,1.20) rectangle (5.78,1.50)}                              % chape de la benne (descend jusqu'à C)
    % --- 2 et 3 : le vérin hydraulique, dessiné dans son repère (origine O, axe x local = u)
    \begin{scope}[shift={(4.2,0.5)}, rotate=150.945]
      \piece{2}{(0,0) circle (0.19)}                                          % œil du pied de vérin
      \piece{2}{(0.10,-0.20) rectangle (0.28,0.20)}                           % fond du corps
      \piece{2}{(0.28,0.13) rectangle (1.50,0.20)}                            % paroi du corps (coupée)
      \piece{2}{(0.28,-0.20) rectangle (1.50,-0.13)}
      \piecevue{3}{(0.32,-0.13) rectangle (0.47,0.13)}                        % piston (non coupé)
      \piecevue{3}{(0.47,-0.065) rectangle (2.0591,0.065)}                    % tige
      \piecevue{3}{(2.0591,0) circle (0.16)}                                  % œil de la tige
      \draw[dash pattern=on 5pt off 2pt on 1pt off 2pt, line width=0.35pt] (-0.35,0) -- (2.35,0);
      \piecevue{5}{(0,0) circle (0.075)}                                      % axe du pied de vérin (en O)
      \piecevue{6}{(2.0591,0) circle (0.075)}                                 % axe tige/benne (en P)
    \end{scope}
    % --- axe de basculement (en C)
    \piecevue{4}{(5.6,1.2) circle (0.13)}
    % --- roues (non coupées), solidaires du châssis : camion freiné
    \piecevue{0}{(1.05,0.42) circle (0.42)}
    \piecevue{0}{(1.05,0.42) circle (0.15)}
    \piecevue{0}{(4.85,0.42) circle (0.42)}
    \piecevue{0}{(4.85,0.42) circle (0.15)}
    % --- points caractéristiques
    \foreach \p/\n/\pos in {C/{(5.6,1.2)}/{}, O/{(4.2,0.5)}/{}, P/{(2.4,1.5)}/{}} {
      \draw[blue, line width=0.5pt] \n ++(-0.13,0) -- ++(0.26,0) \n ++(0,-0.13) -- ++(0,0.26);
    }
    \node[blue, font=\small, inner sep=1pt] at (5.28,1.02) {C};
    \node[blue, font=\small, inner sep=1pt] at (4.40,0.75) {O};
    \node[blue, font=\small, inner sep=1pt] at (2.15,1.38) {P};
    \bnlabels
  \end{scope}}

\begin{document}
\begin{tikzpicture}[x=1cm,y=1cm]
% E0 = {0, 4, 5} châssis (noir) ; E1 = {1, 6} benne ; E2 = {2} corps ; E3 = {3} tige
\newcommand\piececol[3]{\draw[line width=0.6pt, line join=round, fill=#1!#2, draw=#1] #3;}
\newcommand\piece[2]{\ifcase#1\relax
  \piececol{black}{12}{#2}%
  \or \piececol{solideB}{45}{#2}%
  \or \piececol{solideE}{45}{#2}%
  \or \piececol{solideA}{40}{#2}%
  \or \piececol{black}{12}{#2}%
  \or \piececol{black}{12}{#2}%
  \or \piececol{solideB}{45}{#2}%
  \fi}
\newcommand\piecevue[2]{\piece{#1}{#2}}
\newcommand\bnlabels{%
  \node[fill=black, text=white, font=\small\bfseries, inner sep=2pt, rounded corners=1pt] at (1.55,-0.60) {E0 = \{0, 4, 5\}};
  \draw[black, line width=0.5pt] (1.35,-0.40) -- (1.15,0.20);
  \node[fill=solideB, text=white, font=\small\bfseries, inner sep=2pt, rounded corners=1pt] at (3.60,2.85) {E1 = \{1, 6\}};
  \draw[solideB, line width=0.5pt] (3.60,2.65) -- (3.60,1.70);
  % E2 et E3 ne contiennent qu'une pièce : étiquette courte, seule qui tienne entre châssis et benne
  \node[fill=solideE, text=white, font=\small\bfseries, inner sep=2pt, rounded corners=1pt] at (2.05,0.72) {E2};
  \draw[solideE, line width=0.5pt] (2.32,0.72) -- (3.19,0.84);
  \node[fill=solideA, text=white, font=\small\bfseries, inner sep=2pt, rounded corners=1pt] at (2.05,1.06) {E3};
  \draw[solideA, line width=0.5pt] (2.32,1.06) -- (2.87,1.24);
  \node[solideE, font=\footnotesize\bfseries, anchor=west] at (3.30,-0.60) {E2 = \{2\}};
  \node[solideA, font=\footnotesize\bfseries, anchor=west] at (4.80,-0.60) {E3 = \{3\}};
}
\benne
\repere{(0.15,-1.25)}{x}{y}{1}
\end{tikzpicture}
\end{document}
